\documentclass[10pt,a4paper]{amsart}
\usepackage[margin=19mm]{geometry}
\usepackage{amsmath,amssymb,mathtools}
\usepackage[scr=rsfs,cal=euler]{mathalfa}
\usepackage{xfrac}
\usepackage[unicode,hidelinks,bookmarks=false]{hyperref}
\newcommand{\ie}{i.e.\ }
\newcommand{\cf}{cf.\ }
\newcommand{\resp}{respectively\ }
\newcommand{\Subsection}{\subsection}
\providecommand{\nobreakdash}{\nobreak\hbox{-}}
\setlength{\emergencystretch}{2em}
\multlinegap0pt
\usepackage{xcolor}
\usepackage{amssymb}
\usepackage{mathtools}
\usepackage{tikz}
\usepackage[shortlabels]{enumitem}
\usepackage{tcolorbox}
\newtheorem{lemma}{Lemma}
\newcommand{\AKO}{{A_{\tilde K,0}}}
\newcommand{\IL}{\mathbb{L}}
\newcommand{\IP}{\mathbb{P}}
	\newcommand{\Kt}{\widetilde{K}}
\newcommand{\const}[1]{C_{\eqref{#1}}}
\newcommand{\eps}{\varepsilon}
	\newcommand{\hKot}{h_{1,K}}
	\newcommand{\hKtt}{h_{2,K}}
\newcommand*{\norm}[1]{\left\|#1\right\|}
\title{Embedded Trefftz DG method for reaction--diffusion problems on anisotropic meshes}
\author{Sergio G\'omez, Chiara Perinati, Paul Stocker, Igor Voulis}
\date{Source: arXiv:2606.03845v1 (CC BY 4.0)}
\begin{document}
\maketitle

\noindent\emph{Source and licence.} Embedded Trefftz DG method for reaction--diffusion problems on anisotropic meshes. Version arXiv:2606.03845v1 (CC BY 4.0), distributed by arXiv under the Creative Commons Attribution 4.0 International licence, http://creativecommons.org/licenses/by/4.0/, as recorded in the arXiv licence metadata for this exact version and shown on the abstract page https://arxiv.org/abs/2606.03845v1. Full licence notice: \texttt{licenses/arxiv-2606.03845-CC-BY-4.0.txt}.\par\smallskip

\noindent\textit{Editorial scope.} Counted unit: the article's lemma lem:localcoercivity (source0.tex lines 979--994). The article's complete original proof of it is delivered with it (source0.tex lines 995--1188, one \texttt{proof} environment of 194 lines). No mathematics is written, repaired or re-derived: the statement and the proof are the article's own lines, in the article's own notation, and no retained line is substituted or re-flowed.  Retained uncounted as the source-local context the proof needs: lines 907--909; lines 933--940.  Nothing was omitted from any retained span, so the package carries no omission record.  No retained line contains a citation, so the wrapper carries no bibliography and no bibliography entry.  No retained line contains a figure environment, an image-inclusion command or drawing code, so no image is delivered and the package declares no asset dependency.  Editorial formatting only, in addition: an AMS \texttt{amsart} preamble that restates the article's own declarations and macros used by the retained lines, namely the environments \texttt{lemma} on separate counters, together with the standard packages those lines need.  The retained environments print in this document as Lemma 1 in order of appearance, whereas the article prints the same items with the numbering of its own full document.  The complete native source of the article as distributed in the arXiv e-print for this exact version is bundled unchanged as \texttt{sources/202008/source0.tex}.
	\begin{equation}\label{ainv}
	    \|\partial_2 \tilde v\|_{\Kt} \leq \const{ainv} \frac{\hKot}{\hKtt} \|\partial_1 \tilde v\|_{\Kt} \quad \forall \tilde v \in \IL_h(\Kt).
	\end{equation}

	\begin{lemma}\label{lem:noeigenpoly}
	Given a polynomial degree $p\in\mathbb{N}$, 
    we have 
	\begin{equation}\label{eq:noeigenpoly}
		|(\phi'', \phi)_{[0,1]} | \leq (1-C_p)\,\norm{\phi''}_{[0,1]} \,\norm{\phi}_{[0,1]} \qquad \forall \phi \in \IP^p([0, 1]), 
		\end{equation}
		for some constant $0<C_p<1$ depending only on $p$. 
		\end{lemma}

	\begin{lemma}[Stability estimate for~$A_{\Kt, 0}$]\label{lem:localcoercivity}
		There exists a constant $c_{A,0}>0$ such that
	\begin{equation*}
		\|\AKO  \tilde v \|_{\IL_h(\Kt)'} 
	    \geq c_{A,0}
	    \Big( \norm{\tilde v}_{\Kt}^2 + \eps^2 \norm{\partial_1 \tilde v}_{\Kt}^2 + \Big(\const{ainv}\frac{\hKot}{\hKtt}\Big)^{-2} \eps^2 \norm{\partial_2 \tilde v}_{\Kt}^2 \Big)^{\frac{1}{2}}
	    \quad \forall \tilde v\in \IL_h(\Kt),
	\end{equation*}
    and, in particular,
    it holds 
	\begin{equation*}
		\|\AKO  \tilde v \|_{\IL_h(\Kt)'} 
        \geq c_{A,0} \norm{\tilde v}_{\mathbb{L}_h(\Kt)}
	    \quad \forall \tilde v\in \IL_h(\Kt).
	\end{equation*}
	\end{lemma}

	\begin{proof}
	For \(p=1\), the space \(\IL_h(\Kt)\) is trivial, so there is nothing to prove. Hence assume \(p\ge2\).
	
	Let \(\{\psi_i\}_{i=1}^{p-1}\subset \IP^p(\Kt_1)\cap H_0^1(\Kt_1)\) be a basis that is orthonormal in $L^2(\Kt_1)$ and orthogonal with respect to the inner product
$
	(\cdot,\cdot)_{\Kt_1}
	+\eps^2(\partial_1\cdot,\partial_1\cdot)_{\Kt_1}
$,  i.e., such that
	\begin{equation}\label{eq:d1ortho}
	(\psi_i,\psi_j)_{\Kt_1}
	+\eps^2(\partial_1\psi_i,\partial_1\psi_j)_{\Kt_1}
	=\lambda_i\delta_{ij},
	\qquad i,j=1,\dots,p-1.
	\end{equation}
	In particular, for \(i\neq j\),
	$
	(\psi_i,\psi_j)_{\Kt_1}=0$, and
	$(\partial_1\psi_i,\partial_1\psi_j)_{\Kt_1}=0$, and there exist $\widehat\lambda_i>0$, depending only on $p$, such that
	\begin{equation*}
    \lambda_i=1+\eps^2\hKot^{-2}\widehat\lambda_i.
	\end{equation*}
	
	Let \(\tilde v\in \IL_h(\Kt)\) and write
	\begin{equation*}
	\tilde v(\widetilde x_1,\widetilde x_2)=\sum_{i=1}^{p-1}\psi_i(\widetilde x_1)\tilde v_i(\widetilde x_2),
	\quad \text{for }\tilde v_i\in \IP^p(\Kt_2),  i = 1, \ldots, p-1.
	\end{equation*}
	
	The functions $\{\psi_j  \tilde v_j\}_{j=1}^{p-1}$ 
    are orthogonal on $\Kt$ with respect to several operators. More precisely, for all $i \neq j$, we have
	\begin{equation}\label{eq:psiortho}
	\begin{cases}
	( \psi_i  \tilde v_i, \psi_j  \tilde v_j )_{\Kt} = 0,\\
	(\psi_i  \partial_2^2 \tilde v_i, \psi_j  \tilde v_j )_{\Kt} = 0,\\
	( \partial_1^2 \psi_i  \tilde v_i, \psi_j  \tilde v_j )_{\Kt} = 0.
	\end{cases}
	\end{equation}
	The first two equalities follow from the orthogonality of $\{\psi_i\}_{i=1}^{p-1}$ in $L^2(\Kt_1)$ and the tensor-product structure of~$\Kt$. For the last one, using integration by parts in the $\widetilde x_1$-variable and the fact that $\psi_j=0$ on $\partial\Kt_1$, we obtain
	\begin{equation*}
	(\partial_1^2\psi_i\,\tilde v_i,\psi_j\,\tilde v_j)_{\Kt}
	=
	-(\partial_1\psi_i,\partial_1\psi_j)_{\Kt_1}
	(\tilde v_i,\tilde v_j)_{\Kt_2}
	=0,
	\end{equation*}
	for \(i\neq j\), where we used \eqref{eq:d1ortho}.
	Using \eqref{eq:psiortho} and \eqref{eq:d1ortho}, we obtain the following expression for the norm of $\tilde v$:
	\begin{equation}\label{eq:norm-psiortho}
		\begin{split}
			\|\tilde v\|_{\mathbb L_h(\Kt)}^2
			=
			\|\tilde v\|_{\Kt}^2
			+\eps^2\|\partial_1\tilde v\|_{\Kt}^2
			&=
			\sum_{i=1}^{p-1}
			\left(
			\|\tilde v_i\|_{\Kt_2}^2
			+\eps^2\|\partial_1\psi_i\|_{\Kt_1}^2
			\|\tilde v_i\|_{\Kt_2}^2
			\right)
			\\
			&=
			\sum_{i=1}^{p-1}\lambda_i\|\tilde v_i\|_{\Kt_2}^2
			=
			\sum_{i=1}^{p-1}\lambda_i\|\psi_i\tilde v_i\|_{\Kt}^2 .
		\end{split}
	\end{equation}
	Now let \(\tilde u\in \IL_h(\Kt)\) and write 
	\begin{equation*}
	\tilde u(\widetilde x_1,\widetilde x_2)
	=
	\sum_{i=1}^{p-1}\psi_i(\widetilde x_1)\tilde u_i(\widetilde x_2),
	\quad  \text{for } \tilde{u}_i\in \IP^p(\Kt_2), \ i = 1, \ldots, p-1.
	\end{equation*}
	For any $i\in\{1,\dots,p-1\}$, set
	\begin{equation*}
	u_i:=\psi_i\tilde u_i,
	\qquad
	v_i:=u_i-\frac{\eps^2}{\lambda_i}\partial_2^2u_i
	=
	\psi_i\Big(\tilde u_i-\frac{\eps^2}{\lambda_i}\partial_2^2\tilde u_i\Big)\in \IL_h(\Kt).
	\end{equation*}
    
	By \eqref{eq:d1ortho}, these choices satisfy
	$ (u_i-\eps^2\partial_1^2u_i,v_i)_{\Kt} = \lambda_i(u_i,v_i)_{\Kt}$.
	Therefore,
	\begin{align*}
	(\AKO u_i,v_i)_{\Kt}
	&=
	(u_i-\eps^2\partial_1^2u_i-\eps^2\partial_2^2u_i,v_i)_{\Kt}
	\\
	&=
	\lambda_i(u_i,v_i)_{\Kt}
	-(\eps^2\partial_2^2u_i,v_i)_{\Kt}
	\\
	&=
	\lambda_i\|u_i\|_{\Kt}^2
	-2(u_i,\eps^2\partial_2^2u_i)_{\Kt}
	+\frac{1}{\lambda_i}\|\eps^2\partial_2^2u_i\|_{\Kt}^2,
	\end{align*}
	where we used integration by parts, the fact that $\psi_i$ vanishes on the boundary $\partial \Kt_1$, and the 
    orthogonality relations~\eqref{eq:psiortho}.
	
	Using Lemma~\ref{lem:noeigenpoly} in the \(\widetilde x_2\)-variable and the weighted Young inequality, we get
	\begin{align*}
	2\big|(u_i,\eps^2\partial_2^2u_i)_{\Kt}\big|
	&\le
	2(1-C_p)\|u_i\|_{\Kt}
	\|\eps^2\partial_2^2u_i\|_{\Kt}
	\\
	&\le
	(1-C_p)
	\Big(
	\lambda_i\|u_i\|_{\Kt}^2
	+\frac1{\lambda_i}\|\eps^2\partial_2^2u_i\|_{\Kt}^2
	\Big),
	\end{align*}
	and we infer
	\begin{equation}\label{eq:mode-coercive}
	(\AKO u_i,v_i)_{\Kt}
	\ge
	C_p
	\Big(
	\lambda_i\|u_i\|_{\Kt}^2
	+\frac1{\lambda_i}\|\eps^2\partial_2^2u_i\|_{\Kt}^2
	\Big).
	\end{equation}
	Moreover, since \eqref{eq:d1ortho} also gives $ \|v_i\|_{\mathbb L_h(\Kt)}^2 = \lambda_i\|v_i\|_{\Kt}^2$,
	 then by \eqref{eq:mode-coercive} and the triangle inequality, 
	\begin{align*}
	\frac{(\AKO u_i,v_i)_{\Kt}}{\|v_i\|_{\mathbb L_h(\Kt)}}
	&\ge
	C_p\,
	\frac{
	\lambda_i\|u_i\|_{\Kt}^2
	+\lambda_i^{-1}\|\eps^2\partial_2^2u_i\|_{\Kt}^2
	}{
	\lambda_i^{1/2}\|v_i\|_{\Kt}
	}
	\\
	&\ge
	C_p\,
	\frac{
	\lambda_i\|u_i\|_{\Kt}^2
	+\lambda_i^{-1}\|\eps^2\partial_2^2u_i\|_{\Kt}^2
	}{
	\lambda_i^{1/2}\|u_i\|_{\Kt}
	+\lambda_i^{-1/2}\|\eps^2\partial_2^2u_i\|_{\Kt}
	}
	\ge
	\frac{C_p}{2}\,\lambda_i^{1/2}\|u_i\|_{\Kt}.
	\end{align*}
	For every $i=1,\dots,p-1$, using again the orthogonality relations \eqref{eq:psiortho}, we have
	$
	(\AKO \tilde u,v_i)_{\Kt}
	=
	(\AKO u_i,v_i)_{\Kt}
	$
so that
	\begin{equation*}
	\|\AKO\tilde u\|_{\mathbb L_h(\Kt)'}
	=
	\sup_{q\in \mathbb L_h(\Kt)}
	\frac{(\AKO\tilde u,q)_{\Kt}}{\|q\|_{\mathbb L_h(\Kt)}}
	\ge
	\frac{(\AKO\tilde u,v_i)_{\Kt}}{\|v_i\|_{\mathbb L_h(\Kt)}}
	\ge
	\frac{C_p}{2}\lambda_i^{1/2}\|u_i\|_{\Kt} .
	\end{equation*}
	Since this holds for every \(i\), taking the maximum over \(i\) and using
	\(
	\max_{1\le i\le p-1} a_i \ge \frac1{p-1}\sum_{i=1}^{p-1} a_i
	\)
	for nonnegative numbers \(a_i\), we get 
	\begin{align*}
	\|\AKO\tilde u\|_{\mathbb L_h(\Kt)'}^2
	&\ge
	\frac{C_p^2}{4}
	\max_{1\le i\le p-1}\lambda_i\|u_i\|_{\Kt}^2
	\ge
	\frac{C_p^2}{4(p-1)} \sum_{i=1}^{p-1}\lambda_i\|u_i\|_{\Kt}^2.
	\end{align*}
	Using the anisotropic inverse estimate \eqref{ainv} and 
    identity~\eqref{eq:norm-psiortho}, we finally obtain
	\begin{align*}
	\|\AKO\tilde u\|_{\mathbb L_h(\Kt)'}^2
	&\gtrsim
	\|\tilde u\|_{\Kt}^2
	+\eps^2\|\partial_1\tilde u\|_{\Kt}^2
	+\Big(\const{ainv}\frac{\hKot}{\hKtt}\Big)^{-2}\eps^2\|\partial_2\tilde u\|_{\Kt}^2
	.
	\end{align*}
	This proves the claim.
	\end{proof}

\end{document}
